\chapter{Estudo de Conversão}

\section{Análise de Aeronaves de Referência}

O desenvolvimento do conceito de operação em missão dupla exige uma arquitetura de cabine que minimize o tempo de pátio durante o processo de reconfiguração. Na análise de projetos consolidados, observam-se filosofias distintas para a solução deste problema.

Aeronaves regionais civis tradicionais, como o ATR-72, quando adaptadas para operações de carga ou conversões rápidas, dependem da extração de assentos e inserção de cargas através de portas laterais. Na versão \textit{Quick Change} ou nos layouts convertíveis do ATR-72, o processo logístico exige a movimentação de componentes unitários através de aberturas dispostas lateralmente na fuselagem. O processo oficial de conversão de cabine leva em torno de uma hora no ATR-72, sendo frequentemente planejado como um procedimento para as janelas noturnas.

Em geral, no procedimento de conversão, a equipe de solo inicia com a desmontagem manual de elementos de cabine, como os painéis de divisão, as redes de contenção e o assento dianteiro do comissário. Após o destravamento mecânico das fileiras de assentos, que são instaladas de modo convencional nos trilhos do piso, estas fileiras são transladadas manualmente ao longo do corredor central em direção às saídas operacionais. A remoção das poltronas para o exterior e a inserção da carga paletizada ocorre através de uma porta de carga alargada (conhecida como \textit{Large Cargo Door}). Esta porta dispõe de uma soleira rebaixada a uma certa altura do solo para facilitar o recebimento dos contêineres otimizados.  

Embora o sistema QC atenda às normas e permita ao ATR acomodar ULDs no main deck, a cinemática de conversão por portas laterais exigem que os assentos e paletes extraídos precisem ser rotacionados em 90 graus no interior da aeronave antes de passarem pelo vão da porta, o que demanda maior intervenção braçal e expõe o revestimento interno a danos.

Tomando como outro exemplo a aeronave CASA C-295, originada como um cargueiro tático, apresenta uma arquitetura intrinsecamente modular, resolvendo o gargalo direcional através de uma rampa traseira e um piso plano provido de um sistema contínuo de roletes e pontos de fixação compatíveis com paletes. A transição total da cabine de carga para uma configuração completa e segura de passageiros ou transporte aeromédico pode ser finalizada em menos de 2 horas.

Em geral, os assentos de passageiros são fixados previamente de forma integral sobre plataformas padronizadas de assoalho (como o sistema multibase), juntamente com equipamentos auxiliares de oxigênio e segurança. A montagem estrutural de cada palete fora da aeronave é rápida, exigindo apenas dois operadores e cerca de 20 minutos. Para a inserção, a rampa traseira principal atua como a interface direta entre o nível de carregamento e a cabine, podendo operar nivelada com equipamentos de rampa de solo. O translado logístico dos grandes paletes modulares é realizado com o auxílio de um guincho mecânico no interior da cabine. Finalmente, os paletes contendo os assentos são fixados de modo consolidado no piso, sem a necessidade de acoplar as fileiras individualmente.

Essa solução longitudinal ininterrupta suprime a extração avulsa e contorna as limitações de rotação das portas laterais. Contudo, é importante ponderar as penalidades de projeto inerentes a esta arquitetura. A necessidade de uma rampa traseira pressurizável e de um assoalho com trilhos estruturais eleva substancialmente o peso vazio operacional (OEW) da aeronave, além de o formato elevado da cauda impor um incremento considerável de arrasto aerodinâmico de base.

\section*{Escolha e Justificativa da Arquitetura de Carregamento}

Para este projeto, optou-se pela implementação de uma porta de carga lateral (\textit{Large Cargo Door} - LCD) em detrimento de uma rampa de carregamento traseira. Do ponto de vista da integração estrutural, a porta lateral é significativamente mais simples de projetar e manufaturar, pois preserva a geometria cilíndrica convencional da fuselagem. A adoção de uma porta traseira exigiria o projeto de um cone de cauda elevado, o que não apenas incrementa o arrasto aerodinâmico de base, mas também pode alterar de forma indesejada as derivadas de estabilidade e a controlabilidade direcional e longitudinal da aeronave, exigindo um redimensionamento complexo da empenagem.

Adicionalmente, a filosofia de carregamento deve estar estritamente alinhada com o perfil da dupla missão proposta. Rampas traseiras, embora otimizem o tempo de transição, são mandatórias primordialmente para cargas indivisíveis de grande porte ou veículos táticos. Para uma aeronave comercial de transporte regional, cujas cargas consistem essencialmente em paletes padronizados e contêineres ULDs, a porta lateral se mostra a solução mais coerente. Essa escolha simplifica consideravelmente o planejamento aeroportuário, garantindo compatibilidade imediata com a infraestrutura civil existente e permitindo o uso de elevadores de carga e Equipamentos de Apoio no Solo (GSE) convencionais, amplamente disponíveis em aeroportos regionais.

Apesar da adoção da porta lateral, a filosofia de \textit{Quick Change} baseada no deslizamento longitudinal será mantida para maximizar a modularidade da cabine. As fileiras de assentos permanecerão fixadas aos trilhos estruturais do piso. Durante a conversão, os módulos de assentos serão destravados e transladados ao longo do corredor até a seção da porta. Para efetivar a transição física dos assentos e cargas entre o \textit{main deck} e o pátio, a operação dependerá da interface com plataformas elevatórias externas ou rampas móveis posicionadas ao lado da aeronave, compensando a ausência de uma rampa integrada sem penalizar o peso vazio operacional (OEW).

Por fim, é imperativo destacar as mitigações operacionais necessárias inerentes a esta arquitetura. A cinemática de carregamento por porta lateral exige uma inevitável mudança de direção: uma rotação de 90 graus dos paletes e módulos de assentos no interior da aeronave antes da extração. Consequentemente, a área do piso adjacente à porta deverá ser equipada com painéis de transferência omnidirecional e a equipe de \textit{handling} deverá adotar procedimentos rigorosos para evitar impactos e danos estruturais ao revestimento interno durante as manobras logísticas.